\vspace*{-3ex}
\subsection{Event Mining}
\vspace*{-1ex}
\label{sec-learn}

We describe $\MEGA$, $\SVMActive$, and DPF 
that we have recently developed.  
$\MEGA$ is an inductive-based active learning algorithm.
$\SVMActive$ is a margin-based active learning algorithm.
DPF is a perceptual distance function for measuring similarity.
For detailed algorithms and
extensive experimental results, please refer to
\cite{CL01,SVM01,DPF-ACM02}.

\begin{wraptable}{l}{4in}
\footnotesize{
\centerline{
\begin{tabular}{|l|l|l|} \hline
{Filter Name} & {Resolution} & {Representation} \\
\hline \hline
{\it Color Masks} & {Coarse} & {Appearance of culture colors} \\
{\it Color Spread} & {Coarse} & {Spatial concentration of a color} \\
{\it Color Elongation} & {Coarse} & {Shape of a color} \\
\hline
{\it Color Histograms} & {Medium} & {Distribution of colors} \\
{\it Color Average} & {Medium} & {Similarity comparison within} \\
& &  {the same culture color} \\
\hline
{\it Color Variance} & {Fine} & {Similarity comparison within} \\
& & { the same culture color} \\
\hline \hline
\end{tabular}}
}
\caption{Multi-resolution Color Features.}
\label{tab-color-features}
\vspace*{-2ex}
\rule{4in}{0.3mm}
\vspace*{-4ex}
\end{wraptable}	 

Table~\ref{tab-color-features} and Figure~\ref{fig-texture-features}
list the multi-resolution color and texture features that
our work employs for learning image/video query-concepts.
Some individual features\footnote{
This research does not focus on feature-extraction
techniques. We organize useful 
features (e.g., \cite{CBGM97,HCP,HU,HKMZZ,IRM00,MWNS01,PZ,PZM,Rtg98,SC96,Sw91,img10})
in a hierarchical/multi-resolution
manner to assist effective and efficient learning.}
such as color histograms
may be primitive. However,
when organized in a hierarchical/multi-resolution manner,
they can be flexibly chosen by a learning algorithm to model
a subjective concept in different resolutions.
%The design goal of $\MEGA$ and $\SVMActive$ is to learn
%a query concept in this high-dimensional feature space both
%quickly and with a small number of training instances. 



\paragraph{MEGA}
\label{sec-mega}

The Maximizing Expected Generalization Algorithm (MEGA)
models query concepts in $k$-CNF \cite{KLV94},
which can formulate virtually all
practical query concepts\footnote{
$k$-CNF is more expressive than $k$-term DNF,
and it has both polynomial sample complexity
and time complexity \cite{KV94,Mit97}.
}.  MEGA uses
$k$-DNF to bound the sampling space from which to select the
most informative samples to solicit user feedback.

\begin{wrapfigure}{l}{4in}
\begin{center}
\begin{tabular}{c}
\psfig{figure=edfigs/fig-texture-hierarchy.eps,height=3.8cm,width=7cm}
%\psfig{figure=blank.ps,height=3.8cm,width=7cm}
\end{tabular}
\end{center}
\vspace{-.1in}
\caption{Multi-resolution Texture Features.}
\vspace{-.1in}
\label{fig-texture-features}
\vspace*{-2ex}
\rule{4in}{0.3mm}
\vspace*{-1.5ex}
\end{wrapfigure}

\noindent
{\bf Definition 1}:
$k$-{\bf CNF}: For constant $k$, the representation class
$k$-CNF consists of Boolean formulae
of the form $c_1 \wedge \cdots \wedge c_{\theta}$, where
each $c_i$ is a disjunction of at most $k$ literals over
the Boolean variables $x_1, \dots, x_n$. (A Boolean variable
represents an image feature depicted in Table~\ref{tab-color-features} 
and Figure~\ref{fig-texture-features}.)
No prior bound is placed on $\theta$.

\noindent
{\bf Definition 2}:
$k$-{\bf DNF}: For constant $k$, the representation class
$k$-DNF consists of Boolean formulae
of the form $d_1 \vee \cdots \vee d_{\theta}$, where
each $d_i$ is a conjunction of at most $k$ literals over
the Boolean variables $x_1, \dots, x_n$. No prior
bound is placed on $\theta$.

$\MEGA$ initializes
the query concept-space ($QCS$) as a $k$-CNF
and the candidate concept-space ($CCS$) as a $k$-DNF.
The $QCS$ starts as the most specific concept
and the $CCS$ as the most general concept. The target concept
that the learner learns is more general than the initial
$QCS$ and more specific than the initial $CCS$.
The learner learns the $QCS$,
while at the same time refining the $CCS$ to delimit
the boundary of the sampling space. (The
shaded area in Figure~\ref{fig-qccs} shows the sampling space
that is between the $QCS$ and the $CCS$.)

\begin{wrapfigure}{l}{3.5in}
\epsfysize=90pt
\centerline{\epsffile{edfigs/fig-qccs.eps}}
%\centerline{\epsffile{blank.ps}}
\vspace{-.1in}
\caption{MEGA's Sampling Space: $CCS \wedge \overline{QCS}$.}
\vspace{-.1in}
\label{fig-qccs}
\vspace*{-2ex}
\rule{3.5in}{0.3mm}
\vspace*{-1.5ex}
\end{wrapfigure}
                 
Intuitively, at a given stage, we have a $QCS$ and
a $CCS$, representing the boundary of the candidate
concept-space.
To make sure that an example is most useful, we have tested two strategies.
\begin{nlist}
\item {\bf 1}. Bounding the sample space:
Avoid choosing useless unlabeled instances by using
the $CCS$ and $QCS$ to delimit the sampling boundary.
\item {\bf 2}. Maximizing the usefulness of a sample:
Choose an example that
shall remove the maximum expected number of disjunctive
terms from $QCS$. In other words, we choose an example
that can maximize the expected generalization of the concept.
If the example is labeled negative by the user, it can
be useful to remove conjunctive terms in the $CCS$.
\end{nlist}
               
It may appear that
if we pick an example that has more dissimilar
disjunctions (compared to the $QCS$), we may have
a better chance of eliminating more disjunctive terms.
This is, however, not true.
An example must be labeled by the user as positive to be
useful for refining $QCS$. Unfortunately, an example that has
more disjunctions that are dissimilar
to the target concept is less likely to be
labeled positive.  Therefore, there is a tradeoff
between choosing an example that has more contradictory
terms and choosing an example that is more likely
to be labeled positive.
                          
Let $\Psi$ denote the number of disjunctions remaining in the concept.
The number of disjunctions that can be eliminated in the current round
of sampling, denoted as $\psi$, is between zero and $\Psi$.  We can
write the probability of eliminating $\psi$ terms as $P_e(\psi)$.
$P_e(\psi)$ is a monotonically decreasing function of $\psi$.
We attempt to find the $\psi$ that
can eliminate the maximum expected number
of disjunctive terms, given an example.
Our objective function can be written as
\begin{equation}
\psi^{*} = arg_{\psi} \max E(\psi) = \psi \times P_e(\psi).
\label{eq-mega}
\end{equation}
To solve $\psi^{*}$, we must know $P_e(\psi)$, which
can be estimated by the two methods below: probabilistic estimation
and empirical estimation.
(Please refer to \cite{CL01} for detailed discussion.)
We have compared MEGA with the five traditional sampling schemes:
{\it random}, {\it bounded random}, {\it nearest neighbor},
{\it query expansion} and {\it aggressive}.
Our preliminary experimental
results show that when the
number of labeled instances is limited, 
MEGA substantially outperforms traditional
sampling schemes in learning query concepts with high accuracy. 
%Our prototype \cite{demo-Sigmod01} employing $\MEGA$
%can grasp a query concept in three to five
%rounds of feedback in a $200$-feature dimensional space.

\ignore{
\begin{figure}
\begin{center}
\begin{tabular}{llll}
\epsfxsize=80pt
\epsfysize=80pt
%\subfigure[Random]{\epsfbox{fig-mega-random.eps}}
\subfigure[Random]{\epsfbox{blank.ps}}
\hspace{.2in}
\epsfxsize=80pt
\epsfysize=80pt
%\subfigure[B-Random]{\epsfbox{fig-mega-G-random.eps}}
\subfigure[Random]{\epsfbox{blank.ps}}
\hspace{.2in}
\epsfxsize=80pt
\epsfysize=80pt
%\subfigure[Aggressive]{\epsfbox{fig-mega-aggresive.eps}}
\subfigure[Random]{\epsfbox{blank.ps}}
\hspace{.2in}
\epsfxsize=80pt
\epsfysize=80pt
%\subfigure[MEGA]{\epsfbox{fig-mega-mega.eps}} \\
\subfigure[Random]{\epsfbox{blank.ps}}\\
\end{tabular}
\end{center}
\vspace{-.2in}
\caption{Sampling Schemes.}
\vspace{-.1in}         
\label{fig-samplings}
\end{figure}
                    
Figure~\ref{fig-samplings} shows how some of these sampling algorithms work.
The main features of the sampling schemes are given below.
\begin{jlist}
\item {\it Random}: Samples are randomly selected from the bulk of the
domain (Figure~\ref{fig-samplings}(a)).
\item {\it Bounded Random}: Samples are randomly selected from between
$QCS$ and $CCS$  (Figure~\ref{fig-samplings}(b)).
\item {\it Nearest Neighbor}: Samples are selected from the
nearest neighborhood of the center of the positive-labeled instances.
\item {\it Query Expansion}: Samples are selected from the
neighborhood of multiple positive-labeled instances.
\item {\it Aggressive}: Samples are selected from the unlabeled ones
that satisfy the most general concepts in $CCS$
(Figure~\ref{fig-samplings}(c)).
\item {\it MEGA}: Samples are selected from between $QCS$ and $CCS$
and to eliminated the maximum expected number of terms
(Figure~\ref{fig-samplings}(d)).
\end{jlist}
}

\paragraph{$\SVMActive$}
\label{sec-svm}

$\SVMActive$ is 
a method that we have developed
for supporting online concept learning \cite{SVM01}.
$\SVMActive$ combines active learning with
support vector machines (SVMs).
SVMs have strong theoretical foundations and excellent empirical
successes. They have been applied to tasks such as handwritten digit
recognition~\cite{Vapnik:1998}, object
recognition~\cite{Papageorgiou+al:98}, and text
classification~\cite{Joachims:ECML97}.
%We added the active learning component to SVMs for
%selecting the most informative samples to 
%query a user and quickly learns
%a boundary that separates the data objects that satisfy the user's query
%concept from the rest of the dataset.  
%
%\begin{figure}
%\begin{center}
%\psfig{figure=svm2.eps,height=2.4in,angle=90}
%\caption{A simple linear Support Vector Machine.}
%\label{fig:max_margin}
%\end{center}
%\end{figure}

For our purpose of query-concept learning, we
consider SVMs in the binary classification setting. We are
given training data $\{ \vc{x}_1 \ldots \vc{x}_n \}$ that are vectors
in some space $\setX \subseteq \R^d$. We are also given their
labels $\{ y_1 \ldots y_n \}$ where $y_i \in \{ -1,1 \}$. In their
simplest form, SVMs are hyperplanes that separate the training data by
a maximal margin. All vectors lying on
one side of the hyperplane are labeled as $-1$, and all vectors lying
on the other side are labeled as 1. The training instances that lie
closest to the hyperplane are called {\em support vectors}. More
generally, SVMs allow one to project the original training data in
space $\setX$ to a higher dimensional feature space $\setF$ via a
Mercer kernel operator $K$. In other words, we consider the set of
classifiers of the form $f(\vc{x}) = \sum_{i=1}^n
\alpha_i K(\vc{x}_i, \vc{x})$. When $f(\vc{x}) \ge 0$ we classify $\vc{x}$ as
$+1$, otherwise we classify $\vc{x}$ as $-1$.

%\begin{equation}
%\label{eq:kernel_classifier}
%f(\vc{x}) = \sign \left( \sum_{i=1}^n \alpha_i K(\vc{x}_i, \vc{x}) \right).
%\end{equation}

When $K$ satisfies Mercer's condition~\cite{Burges:1998} we can write:
$K(\vc{u}, \vc{v}) = \innerprod{\Phi(\vc{u})}{\Phi(\vc{v})}$ where
$\Phi:\setX \rightarrow \setF$ and ``$\innerprod{}{}$'' denotes an
inner product. We can then rewrite $f$ as:
\begin{equation}
f(\vc{x}) = \innerprod{\vc{w}}{\Phi(\vc{x})}, 
\mbox{ where } \vc{w} = \sum_{i=1}^n \alpha_i \Phi(\vc{x}_i).
\label{eq:kernel_hyperplane}
\end{equation}
Thus, by using $K$ we are implicitly projecting the training data
into a different (often higher dimensional) feature space $\setF$.
The SVM then computes the $\alpha_i$s that correspond to the maximal
margin hyperplane in $\setF$.  By choosing different kernel functions
we can implicitly project the training data from $\setX$ into spaces
$\setF$ for which hyperplanes in $\setF$ correspond to more complex
decision boundaries in the original space $\setX$. 

Intuitively, $\SVMActive$ works by combining the following three ideas:
\begin{enumerate}
\item $\SVMActive$ regards the task of learning a target concept as
one of learning an SVM binary classifier.  An SVM captures the query
concept by separating the relevant images from the irrelevant images
with a hyperplane in a projected space, usually a very
high-dimensional one.  The projected points on one side of the
hyperplane are considered relevant to the query concept and the rest
irrelevant.
\item $\SVMActive$ learns the classifier quickly
via active learning.
The active part of $\SVMActive$ selects the most informative instances
with which to train the SVM classifier. This step ensures fast convergence
to the query concept in a small number of feedback rounds.
\item Once the classifier is trained, $\SVMActive$ returns the top-$k$
most relevant images. These are the $k$ images farthest
from the hyperplane on the query concept side.
\end{enumerate}

Let $V$ denote the uncertain concept space.
The current version of $\SVMActive$ employs
Tong \& Koller's lemma \cite{TK00-b}, which chooses a pool-query
that halves the uncertain space.  
In reality, however, finding the pool-query that
can divide the uncertain space in half can
be difficult and computationally intensive.
Given an unlabeled instance $\vc{x}$ from the pool,
it is not practical to explicitly compute the sizes of the new uncertain
spaces $\setV^{-}$ and $\setV^{+}$ (i.e., the uncertain spaces obtained
when $\vc{x}$ is labeled as $-1$ and $+1$ respectively).  
We have tried a simple method that works as follows:
Learn an SVM on the existing labeled data and choose, as the next
batch of samples, those closest to the 
hyperplane in feature space $\setF$.
Although this simple method works efficiently, 
it has a couple of shortcomings. We discuss 
these shortcomings and how we plan
address them in Section~\ref{sec-future-work}.

%Figure~\ref{fig:simple1}(c) presents an illustration. In the stylized
%picture we have flattened out the surface of the unit weight vector
%hypersphere that appears in Figure~\ref{fig:duality}(a). The white area is
%version space $\setV_i$ which is bounded by solid lines corresponding
%to labeled instances. The five dotted lines represent unlabeled
%instances in the pool. The circle represents the largest radius
%hypersphere that can fit in version space. Note that the edges of the
%circle do not touch the solid lines --- just as the dark sphere
%in~\ref{fig:duality}(b) does not meet the hyperplanes on the surface
%of the larger hypersphere (they meet somewhere under the surface).
%The instance $\vc{b}$ is closest to the SVM $\vc{w}_i$ and so we will
%choose to query $\vc{b}$.

%~\cite{TK00-b} there exist more sophisticated
%approximations that one can perform. However, these methods are
%significantly more computationally intensive. For the task of
%relevance feedback, a fast response time for determining the next 
%data instance
%to present to the user is so critically
%important that these other approximations are not practical.

%The initial theorem of Tong \& Koller assumes one sample
%is chosen for each user iteration.  
%For relevance feedback, we have a need for performing
%multiple pool-queries at the same time. It is not practical to present
%one data object at a time for the user to label --- the user is likely to
%quickly lose patience after a few rounds of pool-querying. To prevent
%this from happening we would like to present the user with multiple
%examples (say, twenty) at each round of pool-querying. Thus, for each
%round, the active learner has to choose not just one example to be
%labeled but, say, twenty.  Theoretically it would be possible to consider
%the size of the resulting uncertain spaces for each possible labeling of
%each possible set of twenty pool-queries but clearly this is
%impractical. Thus instead, for matters of computational efficiency,
%we the simple approach of choosing the pool-queries to be the twenty 
%instances closest to its separating hyperplane.

Once we have performed a number of rounds of querying, we will have
amassed a small set of labeled instances. We learn a final SVM using these
labeled instances and use this SVM as our final classifier. In the 
retrieval task, we wish to return instances in order of decreasing
relevance. This is easy to achieve with an SVM. Recall that an SVM is
a hyperplane in the feature space such that any feature vector lying
on one side of the hyperplane is labeled as ``relevant'' and any
feature vector lying on the opposite side of the hyperplane is labeled
as ``not relevant.'' We have greatest
confidence in those labels when the vectors
are furthest away from the hyperplane. Hence, we can rank all of
the instances in the database in
order of relevance, simply by computing their
signed distance from the hyperplane.


\subsubsection*{\underline{Task: Active Learning with Bayes Point Machines}}

\begin{wrapfigure}{l}{4in}
\label{fig-duality}
\centerline{
\begin{tabular}{ccc}
\psfig{figure=edfigs/vspace1.ps,height=1.4in,angle=270} &
%\psfig{figure=blank.ps,height=1.4in,angle=270} &
\hspace{-0.1in}
\psfig{figure=edfigs/vspace2.ps,height=1.4in,angle=270} \\
%\psfig{figure=blank.ps,height=1.4in,angle=270} \\
(a) & (b)
\end{tabular}}
\footnotesize{
\caption{(a) Version space duality. The surface of the hypersphere
represents unit weight vectors. Each of the two hyperplanes
corresponds to a labeled training instance. Each hyperplane restricts
the area on the hypersphere in which consistent hypotheses can
be situated. Here version space is the surface segment of the hypersphere
closest to the camera. (b) An SVM classifier in version space. The
dark embedded sphere is the largest radius sphere whose center lies in
version space and whose surface does not intersect with the
hyperplanes. The center of the embedded sphere corresponds to the SVM,
its radius is the margin of the SVM in $\setF$, and the training points
corresponding to the hyperplanes that it touches are the support
vectors.}
}
\vspace*{-2ex}
\rule{4in}{0.3mm}
\vspace*{-1.5ex}
\end{wrapfigure}

From a Bayesian perspective Support Vector Machines choose the 
hypothesis corresponding to the largest possible hypersphere 
that can be inscribed in version space, i.e. in the 
space of all consistent hypotheses given a training set \cite{SVM01}.
Given a set of labeled training data and a Mercer kernel $K$, there is
a set of hyperplanes that separate the data in the induced feature
space $\setF$. We call this set of consistent hyperplanes or {\em
hypotheses} the {\em version space}~\cite{Mitchell:AI82}. In other
words, hypothesis $f$ is in version space if for every training
instance $\vc{x}_i$ with label $y_i$ we have that $f(\vc{x}_i) > 0$ if
$y_i = 1$ and $f(\vc{x}_i) < 0$ if $y_i = -1$.
Clearly, by definition, points in $\setW$ correspond to hyperplanes in
$\setF$. The intuition behind the converse is that observing a
training instance $\vc{x}_i$ in feature space restricts the set of
separating hyperplanes to ones that classify $\vc{x}_i$ correctly. In
fact, we can show that the set of allowable points $\vc{w}$ in $\setW$
is restricted to lie on one side of a hyperplane in $\setW$.  More
formally, to show that points in $\setF$ correspond to hyperplanes in
$\setW$, suppose we are given a new training instance $\vc{x}_i$ with
label $y_i$. Then any separating hyperplane must satisfy
$y_i(\innerprod{\vc{w}}{\Phi(\vc{x}_i)}) > 0$. Now, instead of viewing
$\vc{w}$ as the normal vector of a hyperplane in $\setF$, think of
$y_i \Phi(\vc{x}_i)$ as being the normal vector of a hyperplane in
$\setW$. Thus $y_i(\innerprod{\vc{w}}{\Phi(\vc{x}_i)}) =
\innerprod{\vc{w}}{y_i \Phi(\vc{x}_i}) > 0$ defines a half-space in
$\setW$.  Furthermore $\innerprod{\vc{w}}{y_i \Phi(\vc{x}_i}) = 0$
defines a hyperplane in $\setW$ that acts as one of the boundaries to
version space $\setV$. Notice that version space is a connected region
on the surface of a hypersphere in parameter space. 

Figure~$5$(a) presents an example.
SVMs find the center of the largest radius hypersphere whose center
can be placed in version space and whose surface does not intersect
with the hyperplanes corresponding to the labeled instances, as in
Figure~$5$(b). It can be easily shown that the
hyperplanes touched by the maximal radius hypersphere
correspond to the support vectors and that the radius of the
hypersphere is the margin of the SVM.

An alternative and potentially better approach is to construct the 
hypothesis using the whole of version space. This is achieved by using a 
Bayes Point Machine that finds the midpoint of the region of 
intersection of all 
hyperplanes bisecting version space into two halves of equal volume (the 
Bayes point) \cite{Herbrich+al,HG01,HGC01}. It is known that the centre of 
mass of version space 
approximates the Bayes point. The initial work of
Herbrich and Graepel estimates the centre of 
mass by averaging over the trajectory of a billiard 
ball bouncing in version space, sampling version space using a 
Gibbs sampler and approximative methods using the perceptron 
learning algorithm \cite{GH01}.
This method, however, is not efficient.
We plan to investigate new algorithms that can find the 
centre of mass more efficiently. We also plan to
implement active learning with Bayes Point Machines and
compare its performance with $\SVMActive$.


\subsubsection*{\underline{Task $E$: Imbalanced Dataset Classification}}

For many emerging applications,
such as image analysis and gene profiling, where
the training instances of the target class is
heavily outnumbered by the rest of
the training instances,
the class-boundary learned by SVMs can be
severely skewed towards the target class.  As
a result, the false-negative rate in identifying
target objects (e.g., a suspicious event)
can be excessively high
\cite{Kubat97,Japk00,Chaw00,Wei01,Nug02,Chan98,Lin02},
and hence can render the classifier ineffective.
Consider that the number of objects that matches
a query concept is a small fraction of the data.  This
boundary-skew problem can give the learner an even harder
time (in addition to finding rare positive instances)
to compute the class boundary. 

To illustrate this problem, Figure~\ref{fig-checkerboard} depicts 
a 2D checkerboard example. The checkerboard divides a 200*200 
square into four quadrants. The top-left and bottom-right 
quadrants are occupied by negative (majority)
instances, whereas the top-right and bottom-left quadrants contain
positive (minority) instances. The lines between the classes are 
the ``ideal'' boundary that separates the two classes.
In the rest of the paper, we use {\em positive} when referring to
minority instances, and {\em negative} when referring to
majority instances.

{\small
\begin{wrapfigure}{l}{4in}
\vspace{-0.1in}
\begin{center}
\begin{tabular}[t]{c}
\epsfysize=120pt
 \subfigure{\epsfbox{edfigs/cb10.eps}} 
 %\subfigure{\epsfbox{blank.ps}} 
%\epsfysize=160pt
% \subfigure{\epsfbox{cb1000.eps}}
\end{tabular}
\end{center}
\vspace{-0.5in} 
\caption{Checkerboard Experiment.} 
\label{fig-checkerboard}
\vspace*{-2ex}
\rule{4in}{0.3mm}
\vspace*{-1.5ex}
\end{wrapfigure}
}

Figure~$7$ exhibits the boundary distortion between
the two left quadrants in Figure~\ref{fig-checkerboard} under
two different negative/positive training-data ratios.
Figure~$7$(a) shows the SVM class boundary when the
ratio of the number of negative instances (in the quadrant above) to
the number of positive instances (in the quadrant below) is $10:1$. 
Figure~$7$(b) shows the boundary
when the ratio increases to $1,000:1$. 
Compared to the boundary in Figure~$7$(a),
the boundary in
Figure~$7$(b) is much more skewed
towards the positive quadrant, and hence causes
a higher incidence of false negatives.

{\small
\begin{figure}[!h]
\vspace{-0.1in}
\begin{center}
\begin{tabular}[t]{cc}
\epsfysize=110pt
 %\subfigure[10:1]{\epsfbox{de1.eps}} \hspace{.3in}
 \subfigure[10:1]{\epsfbox{blank.ps}} \hspace{.3in}
\epsfysize=110pt
 %\subfigure[1,000:1]{\epsfbox{de1000.eps}}
 \subfigure[1,000:1]{\epsfbox{blank.ps}}
\end{tabular}
\end{center}
\label{fig-boundary}
\vspace{-.2in}
\caption{Boundaries of Different Ratios.}
\end{figure}
}

%To examine the causes, we 
%consider SVMs in the binary
%classification setting. We are given training data
%$\{\textbf{x}_1\cdot\cdot\cdot\textbf{x}_n\}$ that are vectors in
%some space $\chi\subseteq\Re^d$. We are also given their labels
%$\{y_1\cdot\cdot\cdot y_n\}$ where $y_i\in\{-1,1\}$. 
%Suppose the training data is
%linearly separable.  The goal of SVMs is to find the optimal
%hyperplane that separates the training data of
%the two classes by a maximum margin.
%For nonlinearly separable cases, SVMs can project the training
%data onto a higher dimensional feature space via a Mercer kernel
%operator $K$.  In addition, Vapnik's soft-margin theory
%\cite{Vapnik95} introduces slack variables $\xi_i$ to permit
%training error. Given $x_i$, SVMs model class prediction as
%\begin{equation}
% y_i(w\cdot\Phi(x_i)+b)\geq1-\xi_i,~~~\xi_i\geq0,i=1,...,n,
% \end{equation}
%where $w$ is the norm to the hyperplane, $|b|/||w||$ is the
%perpendicular distance from hyperplane to the origin, and $\Phi$
%is an input-to-feature space mapping function. The optimal
%solution of SVMs is formulated by maximizing the Lagrangian
%\begin{equation}
%\label{equ-Lag1}
%L_D=\sum\limits_{i=1}^{n}\alpha_i-\frac{1}{2}\sum\limits_{i,j=1}^{p}\alpha_i\alpha_jy_iy_jK(x_i,x_j),
%\end{equation}
%subject to the following constraints:
% \begin{equation}
%0\leq\alpha_i\leq C~~~~  and~~~~
%\sum\limits_{i=1}^{n}\alpha_i y_i=0,
%\label{equ-constraint}
%\end{equation}
%where $C$ is a constant larger than zero used as penalty for
%soft-margin SVMs. According to the KKT conditions, the value of
%$\alpha_i$ has three ranges \cite{Fle87}:
%
%{\small
%\begin{eqnarray}\label{eq:kkt}
%\alpha_i=0  \Rightarrow  y_i f(x_i)\geq1 & \textmd{and} & \xi_i=0 \\
%0<\alpha_i<C  \Rightarrow  y_i f(x_i)=1  & \textmd{and} & \xi_i=0  \\
%\alpha_i=C  \Rightarrow  y_i f(x_i)\leq1 & \textmd{and} &
%\xi_i\geq0
%\end{eqnarray}
%}
%Once the $\alpha_i$, $i = 1\cdots n$ are solved, the
%class prediction for a testing instance $x$ can be 
%performed by the following equation:
%{\small
%\begin{equation}\label{eq:deci}
%  sgn(f(x)=\sum_{i=1}^n{y_i \alpha_i K(\textbf{x},\textbf{x}_i)}+b).
%\end{equation}
%}
%
%The KKT conditions and Equation~\ref{equ-constraint}
%reveal a couple of important properties of SVMs.
%First, only the training instances that have a non-zero $\alpha_i$ 
%(the support vectors)
%participate in the class decision in Equation~\ref{eq:deci}.
%Second, the constraint in  Equation~\ref{equ-constraint} shows that
%the total weights $\alpha_i$ of the positive support vectors
%must be equal to that of the negative support vectors, or
%{\small
%\begin{equation}\label{eq:dalpha}
%\sum_{i=1}^{np}{\alpha_i}=\sum_{j=1}^{nn}{\alpha_j},
%\end{equation}
%}
%where $np$ and $nn$ denote the number of positive support vectors
%and the number of negative support vectors, respectively.
%
%{\small
%\begin{figure}[!h]
%\vspace{-0.1in}
%\begin{center}
%\begin{tabular}[t]{c}
%\epsfysize=120pt
% \subfigure{\epsfbox{sv_ratio.eps}}
%\end{tabular}
%\end{center}
%\vspace{-.1in} 
%\caption{Support Vectors Ratios} 
%\vspace{-.10in} \label{fig:sv}
%\label{fig-sv-ratios}
%\end{figure}
%}

Using SVM properties along with our preliminary
experimental results,
we observe two potential causes for the boundary skew.
\begin{nlist}
\item {\em 1. The imbalance training-data ratio}.
At the minority side of the boundary, the positive 
training data may not always reside as close to the ``ideal
boundary'' as the negative training data do.  
This can be explained by a toy example.
Suppose we randomly draw $n$ numbers based on 
a uniform distribution between $1$ and $100$.
The larger the value of $n$, the higher the
chance that we draw a number close to $100$,
though the expected mean of the draws is invariant of
the values of $n$.  Thus, the low presence of
the positive training instances makes them locate
farther from the ``ideal boundary'' than 
the negative training instances do.

\item {\em 2. The imbalance support-vector ratio}.
As shown in Figure~$7$, the $\alpha_i$ values of
the minority class (the positive class)
tend to be much larger than those
of the majority class, while the number of positive
support vectors are substantially smaller.  
As a consequence, the nearest
neighborhood of a testing point, especially
when it is near the boundary, is 
likely to be dominated by negative support vectors,
and hence the decision function is
more likely to classify a boundary point negative.
\end{nlist}

We plan to find remedies to the problem by investigating 
the following two methods:
\begin{nlist}
 \item {\em 1. Quasi-bagging}.
  We believe it can be helpful to judiciously downsample the majority
  class.  To remedy the potential
  variance problem caused by sampling, using multiple bags
  can be a good idea.  
 \item {\em 2. Adaptive conformal transformation}.
  Adaptive conformal transformation  gives the neighborhood of
  minority support vectors a higher spatial resolution.
  (The work of \cite{Amari99} proposed a non-adaptive conformal
  transformation method.) We believe this approach can 
  adapt well to data distribution and improve separation
  between the majority and minority classes.
\end{nlist}

